%% Start of main content
% REPLACED 2026-09 by the Dowling Lab fork: the previous document.tex
% was a generic, 2009-era "how Beamer overlays/verbatim/pstricks work"
% tutorial, unrelated to this theme specifically. Testing it directly
% (not just reading it) turned up three independent, pre-existing
% problems -- a broken verbatim block, several commands (\bs) used but
% never defined, and real pstricks drawing commands (\nccurve,
% \pscirclebox, ...) that cannot render under plain pdflatex at all
% (pstricks needs the classic latex+dvips+ps2pdf pipeline; it fails
% with "PSTricks_Not_Configured_For_This_Format" otherwise) -- on top
% of which none of it demonstrated anything specific to this theme.
% Replaced with a short, focused demo of the actual options and
% commands this fork adds, which does compile cleanly under plain
% pdflatex (verified 2026-09, not just written and assumed).

\section{Colors and Background}

\begin{frame}
\frametitle{The independent style options}
Every visual element is its own choice, not a single overall "mode":
\begin{itemize}
  \item \texttt{background} -- white (default) or blue
  \item \texttt{titlebar} -- blue, gold, or none
  \item \texttt{titlecolor}, \texttt{textcolor} -- blue, gold, black, white
  \item \texttt{cornerlogo}, \texttt{monogram}, \texttt{titlepagelogo}
\end{itemize}
Each has a sensible default (see the README), but every combination is
available -- this deck itself is built with
\texttt{\bs usetheme[sectionbar=gold,cornerlogo=blue]\{NotreDame\}}.
(Options are set once for the whole document, in that one
\texttt{\bs usetheme} call -- not per-frame, so a single demo deck
like this one can only show the combination it was actually built
with; \texttt{titlebar=none}, for a bar-free look with the title text
straight on the canvas, is a real option but needs its own separate
build to see live. Also not shown here: \texttt{nav} and
\texttt{fullFooter} both visibly collide with \texttt{cornerlogo}'s
note box, since none of them know about each other -- see the README.)
\end{frame}

\subsection{Highlighting}

\begin{frame}[fragile]
\frametitle{The \bs hl command}
\hl{Bold, colored emphasis} for a phrase, matching a real habit: bold
a few words in a sentence, not the whole thing.
\begin{itemize}
  \item \texttt{\bs hl\{text\}} uses the deck-wide default color (set
    via the \texttt{highlight} option; gold unless changed)
  \item \texttt{\bs hl[gold]\{text\}}, \texttt{\bs hl[blue]\{text\}},
    \texttt{\bs hl[green]\{text\}} always use that color for this one
    call
  \item \texttt{\bs hl[none]\{text\}} just bolds, no color, regardless
    of the deck default
\end{itemize}
For example: \hl[gold]{gold}, \hl[blue]{blue}, \hl[green]{green},
\hl[none]{plain bold}.
\end{frame}

\section{Logos and Marks}

\begin{frame}[fragile]
\frametitle{Corner logo and notes}
\texttt{cornerlogo=<color>} draws a small academic mark in the
bottom-left corner of every frame (this one included) -- \texttt{mark=
primary|secondary} chooses the lockup shape, per Notre Dame's own
guideline that the secondary mark is for when space doesn't allow the
primary one.

\texttt{\bs ndlogonote[align]\{text\}} places a companion note in the
white space beside it: horizontally centered between the logo and the
slide's right edge, vertically centered on the logo's own middle.
[align] defaults to center (a takeaway line); pass [left] for a
reference list.
\ndlogonote{A short takeaway line, centered beside the logo.}
\end{frame}

\begin{frame}[fragile]
\frametitle{The title page mark}
The title page always uses the \textbf{primary} (wide) academic mark,
regardless of \texttt{mark} -- there's always room for it there, per
the same guideline: secondary/monogram are for tight spaces only.
Its color is its own option, \texttt{titlepagelogo}, independent of
\texttt{cornerlogo}.

The small monogram inside the title bar (\texttt{monogram=<color>},
or \texttt{monogram=none}) is the one placement that really is tight
enough to need the compact "ND" letterform instead of the full mark.
\end{frame}
